\documentclass[12pt]{article}

% --------------------- PAQUETES ---------------------
\usepackage[utf8]{inputenc}
\usepackage[T1]{fontenc}
\usepackage[spanish]{babel}
\usepackage{amsmath, amssymb}
\usepackage{graphicx}
\usepackage{float}
\usepackage{xcolor}
\usepackage{booktabs}
\usepackage{siunitx}
\sisetup{output-decimal-marker = {,}} % Para decimales en español

% --------------------- DOCUMENTO ---------------------
\begin{document}

\section*{Pregunta: Test de raíz unitaria y tendencia en series temporales}

\begin{enumerate}

    \item[i)] A partir de la inspección visual de la siguiente serie temporal, ¿considera que el tipo de cambio nominal muestra una tendencia determinista? Justifique su respuesta observando el comportamiento de largo plazo de la serie.

    \begin{center}
    \textbf{Figura 1: Tipo de cambio nominal en Chile (pesos por dólar)}
    \begin{figure}[H]
        \centering
        \includegraphics[width=0.65\linewidth]{certamenes/BCChChart.png}
    \end{figure}
    \end{center}

    \item[ii)] Para evaluar la presencia de raíz unitaria en la serie, utilice el test de Dickey-Fuller aumentado (ADF). Considere el siguiente modelo auxiliar:

    \[
    \Delta y_t = \alpha + \beta t + \phi y_{t-1} + v_t
    \]

    \begin{enumerate}
        \item[a)] ¿Cuál es la hipótesis nula y la hipótesis alternativa del test en este contexto?
        \item[b)] ¿Qué implicancias tiene incluir la tendencia determinista \( \beta t \) en el modelo auxiliar?
        \item[c)] Interprete el modelo bajo la hipótesis alternativa. ¿La serie sería estacionaria alrededor de qué?
    \end{enumerate}

    \item[iii)] A continuación se presentan los resultados de la regresión auxiliar del  test DF bajo dos especificaciones, con y sin tendencia:
   

    \textbf{Modelo sin tendencia (without trend):}
    \begin{table}[H]
    \centering
    \begin{tabular}{lccc}
    \toprule
    Variable & Coef. & Error estándar & Estadístico t \\
    \midrule
    Intercepto & 0,0710 & 0,0354 & 2,006 \\
    $y_{t-1}$  & -0,0182 & 0,0097 & -1,876 \\
    \bottomrule
    \end{tabular}
    \end{table}

    \textbf{Modelo auxiliar con tendencia (\textit{trend}):}
    \begin{table}[H]
    \centering
    \begin{tabular}{lccc}
    \toprule
    Variable & Coef. & Error estándar & Estadístico t \\
    \midrule
    Intercepto & 0,0452 & 0,0289 & 1,563 \\
    Tendencia ($t$) & 0,0021 & 0,0003 & 7,000 \\
    $y_{t-1}$  & -0,0304 & 0,018 & -1,6889 \\
    \bottomrule
    \end{tabular}
    \end{table}

    \begin{enumerate}
        \item[d)] Compare los valores del estadístico t para \( y_{t-1} \) en ambos modelos. ¿Cuál especificación es más adecuada según la evolución de la serie? Justifique.
        \item[e)] A un 1\% de significancia, ¿puede rechazarse la hipótesis nula de raíz unitaria? Use los valores críticos correspondientes al test ADF con constante y con tendencia.
    \end{enumerate}

\end{enumerate}

\end{document}
